\aufzaehlungdrei{Determinan

\mavergleichskettedisp
{\vergleichskette
{ \det \begin{pmatrix} 1+t^2 & t^3 \\ t & 1+t^2 \end{pmatrix}
}
{ =} { 1+2t^2+t^4 -t^4
}
{ =} { 1 +2t^2
}
{ >} {  0
}
{ } {
}
}
{}{}{}
selalu positif. Oleh karena itu, kedua vektor
\mathkor {} {f_1(t)} {dan} {f_2(t)} {}
untuk setiap

\mavergleichskette
{\vergleichskette
{ t
}
{ \in }{  \R
}
{  }{
}
{    }{
}
{   }{
}
}
{}{}{}
membentuk suatu basis.
}{Dari hubungan

\mavergleichskettedisp
{\vergleichskette
{ \begin{pmatrix} f_1 \\f_2 \end{pmatrix}
}
{ =} { \begin{pmatrix} 1+t^2 & t^3 \\ t & 1+t^2 \end{pmatrix} \begin{pmatrix} e_1 \\e_2 \end{pmatrix}
}
{ } {
}
{ } {
}
{ } {
}
}
{}{}{}
menghasilkan

\mavergleichskettedisp
{\vergleichskette
{ \begin{pmatrix} e_1 \\e_2 \end{pmatrix}
}
{ =} { { \frac{ 1 }{ 1+2t^2 } } \begin{pmatrix} 1+t^2 & -t^3 \\ -t & 1+t^2 \end{pmatrix} \begin{pmatrix} f_1 \\f_2 \end{pmatrix}
}
{ } {
}
{ } {
}
{ } {
}
}
{}{}{,}
sehingga

\mavergleichskettedisp
{\vergleichskette
{ e_1
}
{ =} { { \frac{   1+t^2 }{  1+2t^2 } }   f_1 - { \frac{  t^3 }{  1+2t^2 } } f_2
}
{ } {
}
{ } {
}
{ } {
}
}
{}{}{}
dan

\mavergleichskettedisp
{\vergleichskette
{ e_2
}
{ =} { -{ \frac{   t  }{  1+2t^2 } }   f_1 + { \frac{   1+t^2 }{  1+2t^2 } } f_2
}
{ } {
}
{ } {
}
{ } {
}
}
{}{}{.}
}{Untuk koneksi trivial berlaku

\mavergleichskettedisp
{\vergleichskette
{ \nabla_\partial (f)
}
{ =} { \partial f
}
{ } {
}
{ } {
}
{ } {
}
}
{}{}{.}
Dengan demikian,

\mavergleichskettealign
{\vergleichskettealign
{ \nabla_\partial (f_1)
}
{ =} { \nabla_\partial \left(  \left(  1+t^2  \right) e_1 + t^3e_2  \right)
}
{ =} { 2t e_1 + 3t^2 e_2
}
{ =} { 2t  \left(  { \frac{   1+t^2 }{  1+2t^2 } }   f_1 - { \frac{  t^3 }{  1+2t^2 } } f_2   \right)  + 3t^2  \left(  -{ \frac{   t  }{  1+2t^2 } }   f_1 + { \frac{   1+t^2 }{  1+2t^2 } } f_2   \right)
}
{ =} {  \left(   { \frac{   2t+2t^3 }{  1+2t^2 } } - { \frac{   3t^3 }{  1+2t^2 } }  \right) f_1  +  \left(  - { \frac{  t^3 }{  1+2t^2 } }   + { \frac{   3t^2+3t^4 }{  1+2t^2 } }     \right)   f_2
}
}
{
\vergleichskettefortsetzungalign
{ =} { { \frac{   2t - t^3 }{  1+2t^2 } }  f_1  +  { \frac{   3t^2-t^3+3t^4 }{  1+2t^2 } }      f_2
}
{ } {}
{ } {}
{ } {}
}
{}{}
dan

\mavergleichskettealign
{\vergleichskettealign
{ \nabla_\partial (f_2)
}
{ =} { \nabla_\partial \left(  t e_1 + \left(  1+t^2  \right) e_2  \right)
}
{ =} { e_1 + 2t  e_2
}
{ =} { { \frac{   1+t^2 }{  1+2t^2 } } f_1 - { \frac{  t^3 }{  1+2t^2 } } f_2 + 2t \left(  -{ \frac{   t  }{  1+2t^2 } } f_1 + { \frac{   1+t^2 }{  1+2t^2 } } f_2   \right)
}
{ =} { \left( { \frac{   1+t^2 }{  1+2t^2 } }  - { \frac{   2t^2 }{  1+2t^2 } }  \right) f_1  + \left(  - { \frac{   t^3  }{  1+2t^2  } } +  { \frac{   2t+2t^3  }{  1+2t^2 } }  \right) f_2
}
}
{
\vergleichskettefortsetzungalign
{ =} { { \frac{   1 - t^2  }{  1+2t^2 } } f_1 + { \frac{   2t   +t^3 }{  1+2t^2 } }      f_2
}
{ } {}
{ } {}
{ } {}
}
{}{.}
Fungsi-fungsi koefisien tersebut adalah simbol-simbol Christoffel.
}

 [[Kategorie:Latexseite]]
